%% ma: L12-B01-0045
%% lop: 12
%% chuong: 1
%% bai: 1
%% dang: 12.1.8
%% loai: TN4
%% muc: TH
%% dap_an: "C"
%% nguon: "Ảnh giáo viên gửi 05/11/2026 (M11): Kiểm tra giữa HK1 Toán 12 (2026-2027), THPT Hồ Thị Bi (TP.HCM), mã đề 1133, Phần I Câu 21"
%% kien_thuc: "Bài 1 (tính đơn điệu của hàm số bậc ba qua dấu của đạo hàm)"
%% trang_thai: chinh-thuc
%% ghi_chu: "Đề gốc không kèm đáp án; đã giải độc lập. Biến thể thay số của L12-B01-0020 (cùng dạng hàm bậc ba nghịch biến, khác hàm số)."
\begin{ex}
Hàm số $y=-x^3+6x^2+15x-7$ nghịch biến trên khoảng nào dưới đây?
\choice{$(-5;1)$}{$(-1;5)$}{\True $(5;+\infty)$}{$(-\infty;1)$}
\loigiai{$y'=-3x^2+12x+15=-3(x+1)(x-5)$.
$y'<0\Leftrightarrow x<-1$ hoặc $x>5$, nên hàm số nghịch biến trên $(-\infty;-1)$ và $(5;+\infty)$.
Trong các khoảng đã cho chỉ có $(5;+\infty)$ nằm trọn trong khoảng nghịch biến.
Khoảng $(-\infty;1)$ chứa $(-1;1)$ là nơi hàm số đồng biến.}
\end{ex}
